\section{Related Work}
\label{sec:related}

\paragraph{Cache nondeterminism in the engine.}
The llama.cpp server documentation states that prompt caching can make results
nondeterministic because logits differ across batch
sizes~\cite{llamacpp-server}. We do not claim to have discovered that. We
measure when reuse is exact, how the difference varies with architecture,
whether it changes agent actions, and what a mitigation costs.

\paragraph{Cached versus recomputed outputs.}
Chodavarapu and Xu find that KV-cached and cache-free decoding diverge on every
token under FP16 and that FP32 removes the token flips~\cite{chodavarapu-fp16};
Lorup finds the same for a retained cache against a re-prefill in BF16 and
shows that the boundary cache carries the divergent
trajectory~\cite{lorup-stage-replay}. Closest to our work, Patodiya runs a
multi-turn agentic tool-use workload on llama.cpp and vLLM and finds that
cache-hit and recompute paths diverge, increasingly with weight
quantisation~\cite{patodiya-cache-divergence}. That study measures divergence
of trajectories on dense models; we measure changes in the agent's action,
compare hybrid recurrent, sliding-window and dense state, identify the
micro-batch alignment that makes reuse exact, and test a mitigation. Zhang and
Yang reach a changed action by another route: a KV cache kept after an
aborted branch, with the decision tokens held identical, flipped a protected
tool effect in 25 of 63 audited cells, and rebuilding the cache removed
it~\cite{zhang-rollback}. There the cache holds content that the transcript no
longer contains; in our study the cached and recomputed paths see the same
tokens. Concurrently, Gu et al.\ show that approximate, non-prefix KV reuse
makes answers depend on the order of earlier requests~\cite{gu-request-order}.

\paragraph{Numerical nondeterminism.}
He et al.\ trace nondeterminism in LLM serving to kernels whose results depend
on batch size~\cite{tml-nondeterminism}, and Yuan et al.\ show that
floating-point non-associativity changes outputs across batch sizes and
hardware~\cite{yuan-nondeterminism}. vLLM ships a batch-invariant
mode~\cite{vllm-batch-invariant}, and SGLang aligns chunked-prefill
truncation points to a multiple of the split-KV size for
determinism~\cite{sglang-deterministic}. Micro-batch tiling at a reuse boundary is an
instance of the same effect, reached through the prompt cache.

\paragraph{Agent workloads and caching.}
TraceLab characterises coding-agent traces and their prefix-cache hit
rates~\cite{tracelab}; other work improves hit rates and latency for agent
traffic~\cite{vllm-mooncake,policy-runtime}. These studies measure speed. Ours
measures whether the reused state yields the same output.

\paragraph{Hybrid models and prefix caching.}
Prefix caching for hybrid models has been troublesome in several engines. A
llama.cpp issue, since closed, reported full prompt reprocessing on every turn
for a hybrid model~\cite{llamacpp-20225}. vLLM tracks hybrid prefix caching in
an open issue~\cite{vllm-26201}, with reports that caching works only in
blocks of 528 tokens~\cite{vllm-40696} and misses on incremental multimodal
requests~\cite{vllm-43587}. An SGLang roadmap describes its cache's
compatibility with hybrid models as limited~\cite{sglang-21846}. The LMCache
documentation notes that for hybrid models generation ``is not bit-exact
between a cached and a fresh run''~\cite{lmcache-hybrid}, without quantifying
the difference.
